\section*{Chapitre \ref{ch:g9:gravitation} --- La gravitation}

\begin{solution}{exo:g9:gravitation:1}
$F = G \, m_A m_B / d^2$, avec les \omterm{def:g3:measuring-mass:mass}{masses} en \omterm{def:g3:measuring-mass:units}{kilogrammes}, $d$ en
\omterm{def:g2:measuring-length:units}{mètres} entre les centres, $F$ en newtons et
$G = \qty{6.67e-11}{N.m^2/kg^2}$. Avec des mots~: l'attraction croît
avec chacune des \omterm{def:g3:measuring-mass:mass}{masses} et meurt avec le carré de la distance --- et
elle agit également sur les deux partenaires, en sens opposés.
\end{solution}

\begin{solution}{exo:g9:gravitation:2}
Avec exactement \qty{1}{N} --- les attractions de la loi vont toujours
par paires égales et opposées. La pomme accélère visiblement parce que
sa \omterm{def:g3:measuring-mass:mass}{masse} est minuscule~; le même \qty{1}{N} appliqué à
\qty{5.97e24}{kg} de \omterm{def:g4:solar-system:planet}{planète} la remue d'une quantité incommensurable.
\end{solution}

\begin{solution}{exo:g9:gravitation:3}
À $2d$~: le quart de la \omterm{def:g2:pushes-and-pulls:force}{force} à $d$ --- le carré de deux. À $10d$~:
le centième.
\end{solution}

\begin{solution}{exo:g9:gravitation:4}
$F = \num{6.67e-11} \times 80 \times 80 \div 0.25 =
\qty{1.7e-6}{N}$ --- moins de deux millionièmes de newton~: quoi qui
rapproche deux danseurs, ce n'est pas la \omterm{def:g9:gravitation:gravitation}{gravitation}.
\end{solution}

\begin{solution}{exo:g9:gravitation:5}
L'allure à laquelle le boulet tombe, et l'allure à laquelle la surface
de la Terre ronde s'écarte sous lui. Une \omterm{def:g6:sun-earth-moon:axisorbit}{orbite}~: tomber en manquant.
\end{solution}

\begin{solution}{exo:g9:gravitation:6}
À l'altitude de la station, la \omterm{def:g9:gravitation:gravitation}{gravitation} garde presque toute sa
\omterm{def:g2:pushes-and-pulls:force}{force} de surface --- c'est elle qui tient la station sur son cercle.
La station et son équipage sont tous deux en train de \emph{tomber},
ensemble, autour de la Terre~: des compagnons de chute n'appuient sur
rien, et \omterm{def:g1:floating-and-sinking:words}{flotter} est la sensation d'une chute libre partagée.
\end{solution}

\begin{solution}{exo:g9:gravitation:7}
Parce que les cieux traitent en \omterm{def:g3:measuring-mass:mass}{masses} monstrueuses --- \omterm{def:g4:solar-system:planet}{planètes} et
\omterm{def:g4:solar-system:star}{étoiles} multiplient le chétif $G$ en \omterm{def:g2:pushes-and-pulls:force}{forces} puissantes --- et parce
que la \omterm{def:g9:gravitation:gravitation}{gravitation} n'est jamais compensée~: elle ne fait qu'\omterm{def:g1:magnets:magnet}{attirer},
si bien que toute \omterm{def:g3:measuring-mass:mass}{masse} ajoutée creuse encore l'attraction.
\end{solution}

\begin{solution}{exo:g9:gravitation:8}
L'attraction de la \omterm{def:g2:moon-in-the-sky:moon}{Lune} soulève les mers à la fois vers elle (l'eau la
plus proche, la plus fortement attirée) et à l'opposé d'elle (l'eau la
plus éloignée, la moins attirée)~; la Terre en tournant fait passer
chaque côte sous les deux bourrelets chaque jour --- deux marées. Les
plus grandes à la nouvelle et à la \omterm{ex:g2:moon-in-the-sky:shapes}{pleine lune}, quand la main du
Soleil s'aligne sur celle de la \omterm{def:g2:moon-in-the-sky:moon}{Lune}~: les marées de vives-eaux.
\end{solution}

\begin{solution}{exo:g9:gravitation:9}
$F = \num{6.67e-11} \times (\num{7.3e22} \times
\num{5.97e24}) \div (\num{3.8e8})^2 \approx \qty{2.0e20}{N}$ --- un
$2$ suivi de vingt zéros de newtons qui tiennent la \omterm{def:g2:moon-in-the-sky:moon}{Lune} à son bail.
\end{solution}

\begin{solution}{exo:g9:gravitation:10}
Deux réglages tournent l'un contre l'autre~: la \omterm{def:g3:measuring-mass:mass}{masse} de la \omterm{def:g2:moon-in-the-sky:moon}{Lune},
environ quatre-vingts fois plus petite, divise l'attraction par
quatre-vingts~; mais son rayon, près de quatre fois plus petit,
\emph{multiplie} l'attraction de surface par quatre au carré ---
environ quatorze. Quatre-vingts en moins, quatorze en plus~: il reste
près d'un sixième.
\end{solution}

\begin{solution}{exo:g9:gravitation:11}
Le carré inverse amincit la poigne du Soleil avec la distance, et une
poigne plus faible ne peut tenir qu'une chute autour plus lente~: les
\omterm{def:g4:solar-system:planet}{planètes} extérieures sont tenues en laisse lâche et flânent, les
intérieures sont serrées fort et courent --- le sprint de Mercure et
la promenade de 164 ans de Neptune sont la même loi lue à deux
distances.
\end{solution}

\begin{solution}{exo:g9:gravitation:12}
Sa période doit valoir exactement un jour, et il doit tourner vers
l'est --- le sens de rotation propre de la Terre. Satellite et antenne
tournent alors au pas parfait~: la parabole «~immobile~» regarde un
boulet de canon qui tombe autour de la \omterm{def:g4:solar-system:planet}{planète} précisément à l'allure
à laquelle la \omterm{def:g4:solar-system:planet}{planète} tourne sous lui --- immobile aux seuls yeux d'un
observateur qui tourne.
\end{solution}

\begin{solution}{pb:g9:gravitation:1}
\textbf{1.} $F = \num{6.67e-11} \times (0.1 \times
\num{5.97e24}) \div (\num{6.4e6})^2 \approx \qty{0.97}{N}$ --- disons
un newton~: la traction du verger, calculée depuis la constante de
l'Univers.
\textbf{2.} \qty{0.97}{N} de même --- la mutualité~: les attractions
de la loi ne viennent que par paires égales et opposées.
\textbf{3.} Doubler la distance des centres divise l'attraction par
quatre~: environ \qty{0.24}{N}.
\textbf{4.} Chaque miette de la \omterm{def:g4:solar-system:planet}{planète} \omterm{def:g1:magnets:magnet}{attire} la pomme --- montagnes
proches, noyau profond, antipodes lointains --- et la somme de toutes
ces attractions se trouve valoir exactement comme si toute la \omterm{def:g3:measuring-mass:mass}{masse}
siégeait au centre~: un théorème, non une hypothèse, et son auteur
retarda de plusieurs années sa publication jusqu'à savoir le
démontrer.
\textbf{5.} Environ $2000 \times 9.8 \approx \qty{2.0e4}{N}$ --- vingt
mille newtons (la loi avec $d$ égal à un rayon terrestre donne la
même chose).
\textbf{6.} À deux rayons, le carré inverse en laisse le quart~:
environ \qty{4900}{N} pour ses \qty{19600}{N} de \omterm{def:g4:weight-and-mass:weight}{poids} en surface.
\textbf{7.} Oui --- un quart de la \omterm{def:g9:gravitation:gravitation}{gravitation} entière n'est pas un
vide sans pesanteur. Le satellite dépense cette attraction à courber
sa \omterm{def:g6:describing-motion:trajectory}{trajectoire}~: il tombe autour de la Terre en manquant toujours sa
cible, exactement comme l'enseignait le canon de la montagne.
\textbf{8.} «~Tiré assez vite, la chute d'un boulet de canon
s'accorde à la courbure de la Terre, et il tourne pour toujours ---
tombant sans jamais atterrir. Tout satellite est ce boulet, lancé de
côté assez vite pour continuer de manquer le sol.~»
\textbf{9.} $\num{3.8e8} \div \num{6.4e6} \approx 59.4$ rayons
terrestres.
\textbf{10.} Environ $1 \div 59.4^2 \approx 1/3500$ de la \omterm{def:g2:pushes-and-pulls:force}{force} de
surface.
\textbf{11.} Douce parce que la \omterm{def:g9:gravitation:gravitation}{gravitation} à soixante rayons est
diluée trois mille cinq cents fois --- la chute de cinq \omterm{def:g2:measuring-length:units}{mètres} de la
première seconde se réduit à environ un millimètre. Sans fin parce que
ce millimètre de chute s'accorde exactement au millimètre dont la
surface de la Terre s'écarte sous la glissade latérale de la \omterm{def:g2:moon-in-the-sky:moon}{Lune}~: la
chute ne gagne jamais de terrain.
\textbf{12.} Par exemple~: «~La pomme, le satellite et la \omterm{def:g2:moon-in-the-sky:moon}{Lune}
obéissent à une seule phrase~: toute \omterm{def:g3:measuring-mass:mass}{masse} tombe vers toute autre,
selon $G m_A m_B / d^2$ --- la pomme atterrit, le satellite et la \omterm{def:g2:moon-in-the-sky:moon}{Lune}
continuent de manquer leur cible, et la différence n'est qu'une
\omterm{def:g7:motion-average-speed:formula}{vitesse} latérale.~»
\end{solution}
